\documentclass[11pt,a4paper]{article}

%% PACHETE MATEMATICE
\usepackage{amsmath,amssymb,amsfonts}
\usepackage{amsthm}
\usepackage{mathpazo}
\usepackage{eulervm}
\usepackage{enumerate}
\usepackage{color}
\usepackage{graphicx}
\usepackage[all]{xy}
\usepackage{xcolor}
\usepackage{framed}
\usepackage[unicode]{hyperref}
\setcounter{MaxMatrixCols}{20}
\usepackage{multicol}
%\usepackage[romanian]{babel}


%% ASPECT
\linespread{1.05}
\usepackage[T1]{fontenc}
\usepackage[utf8x]{inputenc}
\usepackage[margin=2cm]{geometry}
% \usepackage[color]{showkeys}
\usepackage{anyfontsize}
\usepackage{titlesec}
\titleformat*{\section}{\large\bfseries}

\usepackage{fancyhdr}
\pagestyle{fancy}
\rhead{\small \color{gray} Notițe de Adrian Manea}
\lhead{\small \color{gray} ETTI-AG, 2017---2018, A. Niță \& A. Oprina}


\usepackage[autostyle = false, style = german]{csquotes}
\usepackage[romanian]{babel}
\newcommand\qq{\enquote}

%% COMENZILE MELE
\renewcommand*{\proofname}{Dem.:}
\newcommand\bb{\mathbb}
\newcommand\dr{\mathrm}
\newcommand\kal{\mathcal}
\newcommand\fk{\mathfrak}
\newcommand\op{\oplus}
\newcommand\ot{\otimes}
\newcommand\ds{\displaystyle}
\newcommand\id{\indent}
\newcommand\nid{\noindent}
\newcommand\seq{\subseteq}
\newcommand\sneq{\subsetneq}
\newcommand\speq{\supseteq}
\newcommand\spneq{\supsetneq}
\newcommand\rar{\rightarrow}
\newcommand\Rar{\Rightarrow}
\newcommand\xrar{\xrightarrow}
\newcommand\lar{\leftarrow}
\newcommand\Lar{\Leftarrow}
\newcommand\xlar{\xleftarrow}
\newcommand\lrar{\leftrightarrow}
\newcommand\Lrar{\Leftrightarrow}
\newcommand\ideal{\trianglelefteq}
\newcommand\ai{\text{ a.\^{i}. }}
\newcommand\sk{(\textit{Schi\c{t}\u{a}}) }
\newcommand\rhup{\rightharpoonup}
\newcommand\rhdn{\rightharpoondown}
\newcommand\lhup{\leftharpoonup}
\newcommand\lhdn{\leftharpoondown}
\newcommand\vid{\emptyset}
\newcommand\wh{\widehat}
\newcommand\ol{\overline}
\newcommand\NN{\mathbb{N}}
\newcommand\RR{\mathbb{R}}
\newcommand\ZZ{\mathbb{Z}}
\newcommand\QQ{\mathbb{Q}}
\newcommand\CC{\mathbb{C}}

%% STILURILE MELE
\newtheoremstyle{dotless}{}{}{\itshape}{}{\bfseries}{:}{ }{}

\theoremstyle{dotless}
\newtheorem{theorem}{Teorem\u{a}}[section]
\newtheorem{proposition}{Propozi\c{t}ie}[section]
\newtheorem{lemma}{Lem\u{a}}[section]
\newtheorem{corollary}{Corolar}[section]
\newtheorem{exercise}{Exerci\c{t}iu}

\newtheoremstyle{dotlessdr}{}{}{}{}{\bfseries}{:}{ }{}
\theoremstyle{dotlessdr}
\newtheorem{example}{Exemplu}[section]
\newtheorem{definition}{Defini\c{t}ie}[section]
\newtheorem{remark}{Observa\c{t}ie}[section]




\begin{document}

\begin{center}
  {\large\textbf{Seminar 4 \\ Matricea unei aplicații liniare. Utilizări}}
\end{center}

\vspace{1cm}

\section{Matricea unei aplicații liniare}
\label{sec:mat-ap}

Există o foarte strînsă legătură între matrice și aplicații liniare, definite
pentru orice spațiu vectorial. De aceea, așa cum vom vedea, multe din noțiunile
pe care le știm de la spații vectoriale cu ajutorul aplicațiilor liniare pot fi
reformulate în contextul matricelor.

Punctul de pornire este următorul.

\begin{definition}\label{def:lvw}
  Fie $V$ și $W$ două spații vectoriale peste același corp comutativ $\Bbbk$.

  Definim mulțimea:
  \[
    \kal{L}(V, W) = \{ f : V \to W \mid f \text{ aplicație liniară } \}.
  \]
\end{definition}

Se arată că $\kal{L}(V, W)$ devine chiar un spațiu vectorial peste $\Bbbk$,
cu operațiile:
\begin{align*}
  (f + g)(v) &= f(v) + g(v), \forall v \in V, \forall f, g \in \kal{L}(V, W); \\
  (\alpha f)(v) &= f(\alpha v), \forall \alpha \in \Bbbk, f \in \kal{L}(V, W), v \in V.
\end{align*}

În plus, se poate demonstra următoarea teoremă:

\begin{theorem}\label{thm:mmn-apl}
  În contextul și cu notațiile de mai sus, dacă $\dim V = n, \dim W = m$, atunci
  $\dim \kal{L}(V, W) = m \cdot n$. În particular, $\kal{L}(V, W) \simeq M_{m, n}(\Bbbk)$.
\end{theorem}

Rezultă, așadar, că oricărei aplicații liniare i se poate asocia în mod unic
o matrice. Este vorba despre \emph{matricea aplicației într-o bază}, definită
astfel.

Să luăm un exemplu particular. Considerăm spațiul vectorial real $\RR^3$ și
baza sa canonică, $B = \{ e_1 = (1, 0, 0), e_2 = (0, 1, 0), e_3 = (0, 0, 1) \}$.

Fie aplicația liniară $f : \RR^3 \to \RR^2, f(x, y, z) = (3x - 5y, 2z - x)$.

Matricea aplicației $f$ în baza $B$, notată $M^B(f)$, se definește prin $(f(e_i))^t$.
Adică se calculează $f$ în fiecare element al bazei, iar vectorii obținuți devin
\emph{coloanele} matricei rezultate. Concret:
\[
  \begin{cases}
    f(e_1) &= (3, -1) \\
    f(e_2) &= (-5, 0) \\
    f(e_3) &= (0, 2)
  \end{cases} \Rightarrow %
  M^B(f) =
  \begin{pmatrix}
    3 & -5 & 0 \\
    -1 & 0 & 2
  \end{pmatrix}
\]

Această corespondență este foarte utilă, deoarece, ulterior, în loc să calculăm
valoarea aplicației $f$ într-un vector, se poate \emph{înmulți matricea aplicației %
  cu vectorul respectiv} (pe coloană, așa cum are sens produsul de matrice).

De exemplu, dacă luăm, în cazul de mai sus, vectorul $v = (1, -1, 3) \in \RR^3$,
se arată prin calcul direct că:
\[
  f(v) = (8, 5) = M^B(f) \cdot \begin{pmatrix} 1 \\ -1 \\ 3 \end{pmatrix}.
\]

\textbf{Exercițiu:} Verificați egalitatea de mai sus!

%%%%%%%%%%%%%%%%%%%%%%%%%%%%%%%%%%%%%%%%%%%%%%%%%%%%%%%%%%%%%%%%%%%%%%

\section{Schimbarea matricei la schimbarea bazei}
\label{sec:mat-trecere}

Dacă avem o aplicație liniară $f : V \to W$ și două baze fixate în
cele două spații, atunci matricea lui $f$ în raport cu cele două baze
se obține exprimînd imaginea primei baze în raport cu cea de-a doua.

De exemplu, fie aplicația:
\[
  f : \RR^2 \to \RR^3, \quad f(x_1, x_2) = (3x_1 - 3x_2, x_1 + 2x_2, x_1 - x_2).
\]
Considerăm bazele
\[
  B_1 = \{ v_1 = (1, 1), v_2 = (1, -1) \} \seq \RR^2
\]
și
\[
  B_2 = \{ w_1 = (1, 0, 0), w_2 = (0, 1, 0), w_3 = (0, 0, 1) \} \seq \RR^3.
\]
Atunci, \emph{matricea lui $f$ în raport cu cele două baze} se obține din calculul:
\begin{align*}
  f(v_1) &= (-1, 3, 0) = -1 \cdot w_1 + 3 \cdot w_2 + 0 \cdot w_3 \\
  f(v_2) &= (5, -1, 2) = 5 \cdot w_1 - 1 \cdot w_2 + 2 \cdot w_3.
\end{align*}
Așadar, avem:
\[
  {}^{B_1} M^{B_2} (f) =
  \begin{pmatrix}
    -1 & 5 \\
    3 & -1 \\
    0 & 2
  \end{pmatrix}
\]

Definim acum următorul concept:
\begin{definition}\label{def:matrecere}
  Fie $V$ un $\Bbbk$-spațiu vectorial și $B_1, B_2$ două baze ale sale.

  \emph{Matricea de trecere} de la baza $B_1$ la baza $B_2$ are drept coloane
  coordonatele vectorilor din baza $B_1$ scriși în funcție de baza $B_2$.
\end{definition}

De exemplu, să luăm $V = \RR^3$ și baza canonică $B_1 = \{ e_1, e_2, e_3 \}$,
precum și baza $B_2 = \{ v_1 = (-1, 2, 1), v_2 = (0, -1, -1), v_3 = (1, 1, 1) \}$
(\textbf{Exercițiu:} verificați că $B_2$ este, într-adevăr, bază!).

Atunci avem:
\begin{align*}
  e_1 = (1, 0, 0) &= \alpha_1 v_1 + \alpha_2 v_2 + \alpha_3 v_3 \\
  e_2 = (0, 1, 0) &= \beta_1 v_1 + \beta_2 v_2 + \beta_3 v_3 \\
  e_3 = (0, 0, 1) &= \gamma_1 v_1 + \gamma_2 v_2 + \gamma_3 v_3.
\end{align*}

Rezultă trei sisteme:
\[
  \begin{cases}
    -\alpha_1 + \alpha_3 &= 1 \\
    2 \alpha_1 - \alpha_2 + \alpha_3 &= 0 \\
    \alpha_1 - \alpha_2 + \alpha_3 &= 0
  \end{cases}, \quad
  \begin{cases}
    -\beta_1 + \beta_3 &= 0 \\
    2 \beta_1 - \beta_2 + \beta_3 &= 1 \\
    \beta_1 - \beta_2 + \beta_3 &= 0
  \end{cases}, \quad
  \begin{cases}
    -\gamma_1 + \gamma_3 &= 0 \\
    2 \gamma_1 - \gamma_2 + \gamma_3 &= 0 \\
    \gamma_1 - \gamma_2 + \gamma_3 &= 1
  \end{cases}
\]

Matricea de trecere va fi, atunci:
\[
  {}^{B_1} M^{B_2} =
  \begin{pmatrix}
    \alpha_1 & \beta_1 & \gamma_1 \\
    \alpha_2 & \beta_2 & \gamma_2 \\
    \alpha_3 & \beta_3 & \gamma_3
  \end{pmatrix}
\]

\textbf{Exercițiu:} Finalizați calculele pentru exercițiul de mai sus! (Determinați
explicit coordonatele $\alpha_i, \beta_i, \gamma_i$ și matricea ${}^{B_1} M^{B_2}$.)

Cu aceasta, avem următorul rezultat:
\begin{theorem}\label{thm:sch-trecere}
  Fie $A$ matricea unei aplicații liniare în baza $B$ a unui spațiu vectorial.
  Fie $B'$ o altă bază și $S$ matricea de trecere de la baza $B$ la baza $B'$.

  Atunci matricea aplicației se schimbă după formula:
  \[
    A' = S^{-1} \cdot A \cdot S.
  \]
\end{theorem}

În cazul schimbării bazelor în domeniul și codomeniul unei aplicații
liniare, avem rezultatul:

\begin{theorem}\label{thm:mat-trecere-ambele}
  Fie $V, W$ două $\Bbbk$-spații vectoriale, cu $\dim V = n, \dim W = m$ și
  fie $f \in \kal{L}(V, W)$.

  Fie $B_1 = \{v_i\}$ și $\overline{B}_1 = \{\overline{v}_u\}$ două baze
  în $V$ și $B_2 = \{w_j\}$, respectiv $\overline{B}_2 = \{ \overline{w}_j \}$
  două baze în $W$.

  Fie $A$ matricea de trecere de la $B_1$ la $\overline{B}_1$ și $B$ matricea
  de trecere de la $B_2$ la $\overline{B}_2$.

  Atunci are loc relația:
  \[
    {}^{\overline{B}_1} M^{\overline{B}_2} (f) =
    B^{-1} \cdot {}^{B_1} M^{B_2} (f) \cdot A.
  \]
\end{theorem}

Modificarea coordonatelor unui vector la schimbarea bazei se face conform:

\begin{theorem}\label{thm:sch-coord}
  Fie $x \in V$ un vector care are coordonatele $\pmb{x} = (x_1, \dots, x_n)$ în
  baza $B$ și coordonatele $\pmb{x'} = (x'_1, \dots, x'_n)$ în baza $B'$ ale spațiului
  vectorial $V$.

  Fie $A$ matricea de trecere de la baza $B$ la baza $B'$.

  Atunci legătura între coordonate este dată de:
  \[
    \pmb{x'}^t = A^{-1} \cdot \pmb{x}.
  \]
\end{theorem}

Mai definim:

\begin{definition}\label{def:mat-asemenea}
  Două matrice $P, Q \in M_n(\Bbbk)$ se numesc \emph{asemenea} dacă există o matrice
  inversabilă $A$ astfel încît:
  \[
    Q = A^{-1} \cdot P \cdot A.
  \]

  Notăm relația de asemănare prin ``$\sim$'', deci $P \sim Q$.
\end{definition}

\textbf{Exercițiu:} Arătați că relația de asemănare a matricelor pătratice
de un ordin $n$ dat este o relație de echivalență (reflexivă, simetrică, tranzitivă).

\vspace{.3cm}

Folosind cele de mai sus, rezultă:
\begin{enumerate}[(1)]
\item Două matrice asemenea au același rang.
\item Matricele asociate unei aplicații liniare $f : V \to V$ în două baze
  diferite ale lui $V$ sînt asemenea. Are loc și rezultatul reciproc.
\end{enumerate}

%%%%%%%%%%%%%%%%%%%%%%%%%%%%%%%%%%%%%%%%%%%%%%%%%%%%%%%%%%%%%%%%%%%%%%

\section{Exerciții}
\label{sec:exc}

1. Fie transformarea liniară $f : \RR^3 \to \RR^2$ definită prin:
\[
  f(x_1, x_2, x_3) = (x_1 + 2x_3, x_1 + x_2 - x_3).
\]

\begin{enumerate}[(a)]
\item Verificați teorema rang-defect;
\item Găsiți matricea aplicației $f$ în baza canonică a $\RR^3$;
\item Arătați că $B' = \{ v_1 = (-1, 0, 1), v_2 = (0, -1, 2), v_3 = (-1, 1, -1) \}$
  este bază a lui $\RR^3$;
\item Găsiți matricea lui $f$ în baza $B'$;
\item Găsiți matricea de trecere de la baza canonică $B$ la baza $B'$
\end{enumerate}

\vspace{1cm}

2. Fie transformarea liniară $f \in \kal{L}(\RR^3, \RR^4)$, definită prin:
\[
  f(x_1, x_2, x_3) = (x_1 + x_2, x_1 + x_3, x_2 - x_3, x_1 - x_2 + 2x_3).
\]
\begin{enumerate}[(a)]
\item Verificați teorema rang-defect;
\item Fie bazele:
  \begin{align*}
    B_1 &= \{ u_1 = (1, 1, -1), u_2 = (1, -1, 1), u_3 = (-1, 1, 1) \} \seq \RR^3 \\
    B_2 &= \{ v_1 = (0, 1, 1, 1), v_2 = (1, 0, 1, 1), v_3 = (1, 1, 0, 1), v_4 = (1, 1, 1, 0) \} \seq \RR^4.
  \end{align*}
  Găsiți matricele de trecere de la bazele canonice la $B_1$, respectiv $B_2$;
\item Determinați matricea aplicației $f$ în baza $B_1$, față de baza canonică
  din $\RR^4$;
\item Determinați matricea aplicației $f$ în baza $B_1$, față de baza $B_2$.
\end{enumerate}

\vspace{1cm}

3. Fie $V = \RR[X]_2$ și mulțimea:
\[
  B' = \{ p_1 = 3X^2 + 2, p_2 = X - 5, p_3 = X^2 + 4 \}.
\]
\begin{enumerate}[(a)]
\item Arătați că $B'$ este bază a lui $V$;
\item Determinați matricea de trecere de la baza canonică la baza $B'$;
\item Fie $f : V \to \RR[X]_1$ morfismul de derivare (în raport cu $X$).
  Determinați matricea lui $f$ în baza canonică a lui $V$, față de baza canonică
  alui $\RR[X]_1$, precum și matricea lui $f$ în baza $B'$ a lui $V$
  față de baza $B^{\prime\prime} = \{ q_1 = 7, q_2 = X - 2 \}$ a lui $\RR[X]_1$.
\end{enumerate}

\vspace{1cm}

4*. Fie $V$ spațiul vectorial al funcțiilor reale, generat de baza:
\[
  B = \{1, \cos x, \sin x, \cos 2x, \sin 2x, \cos 3x, \sin 3x \}.
\]
Fie aplicația $\varphi : V \to V$, definită prin:
\[
  (\varphi(f))(x) = 4 \int_0^{2 \pi} \sin^3(x + y) f(y) dy.
\]
\begin{enumerate}[(a)]
\item Arătați că $\varphi$ este aplicație liniară;
\item Determinați $\mathrm{Ker}\varphi$ și $\mathrm{Im}\varphi$;
\item Determinați matricea lui $\varphi$ în baza $B$.
\end{enumerate}

\vspace{1cm}

5*. Fie aplicația urmă $tr : M_n(\RR) \to \RR$, definită prin $tr(A) = \sum a_{ii}$,
unde $A = (a_{ij}) \in M_n(\RR)$.

Arătați că $\mathrm{def}(tr) = n^2 - 1$, unde $\mathrm{def}(tr) = \dim \mathrm{Ker}(tr)$,
se numește \emph{defectul} aplicației $tr$.

%%%%%%%%%%%%%%%%%%%%%%%%%%%%%%%%%%%%%%%%%%%%%%%%%%%%%%%%%%%%%%%%%%%%%%

\section{Exerciții suplimentare}
\label{sec:exc-supl}

6. Fie vectorii $v_1 = (x, 3x, 2), v_2 = (5, -1, x + 4), v_3 = (0, -x, x)$.
\begin{enumerate}[(a)]
\item Determinați $x$ astfel încît mulțimea $B' = \{ v_1, v_2, v_3 \}$ să formeze
  o bază pentru spațiul $\RR^3$.
\item Găsiți matricea de trecere de la baza $B'$ la baza canonică a lui $\RR^3$.
\end{enumerate}

\vspace{1cm}

7. Fie aplicația:
\[
  f : \RR^3 \to \RR^2, \quad f(x, y, z) = (3x - 5y, 2x + z, 5x - 5y + z).
\]
\begin{enumerate}[(a)]
\item Arătați că $f$ este aplicație liniară;
\item Determinați matricea lui $f$ în baza canonică a lui $\RR^3$ și baza canonică
  a lui $\RR^2$;
\item Determinați $\mathrm{Ker}f, \mathrm{Im}f$;
\item Verificați teorema rang-defect și găsiți baze în $\mathrm{Ker}f$ și $\mathrm{Im}f$;
\item Fie mulțimile $B_1 = \{ (-1, 0, 1), (2, 2, 0), (1, 0, 0) \}$ și
  $B_2 = \{ (3, 1), (-2, 2) \}$. Arătați că $B_1$ este bază în $\RR^3$ și
  $B_2$ este bază în $\RR^2$;
\item Găsiți matricele de trecere de la bazele canonice la bazele $B_1$, respectiv $B_2$;
\item Găsiți matricea aplicației $f$ în raport cu bazele $B_1$, respectiv $B_2$.
\end{enumerate}

\vspace{1cm}

8. Fie spațiul $V = \RR^3$ și mulțimile:
\begin{align*}
  S_1 &= \{ (x, y, z) \in \RR^3 \mid 3x + 5y - z = 0 \}; \\
  S_2 &= \{ (x, y, z) \in \RR^3 \mid x = 4z - y \text{ și } z = 2x - y \}.
\end{align*}
\begin{enumerate}[(a)]
\item Arătați că $S_1 \hookrightarrow V$ și $S_2 \hookrightarrow V$;
\item Verificați teorema lui Grassmann pentru $S_1$ și $S_2$, adică:
  \[
    \dim(S_1 + S_2) = \dim S_1 + \dim S_2 - \dim (S_1 \cap S_2).
  \]
\end{enumerate}
\end{document}
